\addcontentsline{toc}{subsection}{Inciso j}

\subsection*{j.}

\textbf{Supón que deseas crear una aplicación para gestión hospitalaria. Considera cada una de las desventajas indicadas en el documento “Purpose of Database Systems”, cuando se administran los datos en un sistema de archivos. Discute la relevancia de cada uno de los puntos indicados, con respecto a la relevancia de cada desafío en el contexto de la gestión de datos de pacientes: historial médico, diagnósticos, tratamientos, citas, acceso a registros médicos, médicos, especialidades, entre otros.}

En la lectura ``Purpose of Database Systems'' \cite{PurposeDatabaseSystems} nos dan ejemplos en un contexto universitario mayormente, de por qué el sistema de archivos no es eficiente y puede traer muchos problemas a cualquier establecimiento que necesite organizar sus datos. En este caso, en el contexto de un hospital y la gestión hospitalaria, procederé a particularizar cada punto para mapear los siete desafíos tratados en la lectura, pero en este contexto específico, centrándome en explicar con ejemplos y notando que cada punto es vital para el hospital.

\begin{enumerate}
    \item \textbf{Redundancia e inconsistencia de datos (Data redundancy and inconsistency):} La redundancia de datos de un paciente es extremadamente peligrosa en el contexto de un hospital. Supongamos que un paciente llega a un hospital por primera vez a la especialidad de urgencias, su malestar es tratado, sus datos son recolectados y se le dirige ahora a cardiología; en esta especialidad se le vuelven a recolectar los datos y se le da un diagnóstico y tratamiento. Días después, el paciente llama para actualizar su número telefónico, pero al no haber centralización de los datos, esta actualización solo se le notifica a la especialidad de urgencias. Esto genera una \textbf{inconsistencia de datos}. Así que, si hay alguna complicación o por alguna razón una de las especialidades necesita del paciente, puede que no tenga el teléfono actualizado del paciente y se tenga que buscar en otras especialidades. Esto hace no solo ineficiente la comunicación, sino que pone en riesgo la salud y la vida de los pacientes.

    \item \textbf{Dificultad en el acceso de datos (Difficulty in accessing data):} Aunque pareciera a primera vista que este punto no afecta directamente la integridad de los pacientes, en realidad sí lo hace de forma indirecta, ya que la dificultad en obtener datos de un hospital ralentiza su funcionamiento y no permite que se atienda a la cantidad de pacientes necesarios, o peor aún, podría atentar contra la salud de los pacientes. Supongamos que el hospital requiere conocer a todos los pacientes a los que se les administró un fármaco, digamos el fármaco X, pero sus datos están almacenados en un sistema de archivos, por lo que no existe un programa que satisfaga esta tarea de manera inmediata. Debido a esto, o se debe buscar manualmente en cada archivo o se debe programar la consulta desde cero; ninguna opción es eficiente y ambas requerirán mucho tiempo. Esto puede ser muy peligroso si la razón por la que se necesita conocer a qué pacientes se les administró el fármaco X es porque este podría provocar reacciones alérgicas no identificadas antes. Esta es una de las diversas adversidades que pueden surgir por una mala elección del método de almacenamiento de datos.

    \item \textbf{Aislamiento de datos (Data isolation):} Aquí podemos ver un mapeo casi directo de lo dicho en la lectura ``Purpose of Database Systems'' \cite{PurposeDatabaseSystems} en la parte de aislamiento de datos. Si cada especialidad del hospital está programada por distintos equipos, al querer escribir un nuevo programa (como el del punto anterior para obtener la lista de los pacientes a los que se les administró el fármaco X), implementar dicho programa en las diversas especialidades sería un problema aún peor de lo que se pensaba.

    \item \textbf{Problemas de integridad (Integrity problems):} Este punto parece similar al anterior, pero en este caso nos centraremos en la constante actualización del sistema, en particular para implementar restricciones de consistencia. Por ejemplo, si el hospital adquiere un nuevo medicamento, este tendrá la restricción de dosis máxima segura, digamos 50 mg. Como este fármaco se utilizará en prácticamente todas las especialidades del hospital, se necesitará implementar esta restricción; sin embargo, nos enfrentamos al problema de que cada especialidad tiene su propia implementación independiente, por lo que su aplicación coordinada será sumamente complicada y se podría vulnerar la integridad de los datos.

    \item \textbf{Problemas de atomicidad (Atomicity problems):} Este es uno de los problemas más críticos que debe tener en cuenta el hospital, ya que debe tener un plan de acción ante fallas que, como en cualquier otro sistema computacional, pueden ocurrir. El concepto plantea que las transacciones o movimientos deben ser atómicos, es decir, que se realiza toda la transacción o no se hace nada de ella. Supongamos que se le va a asignar un quirófano y un médico a una cirugía de apéndice, pero en ese momento ocurre una falla eléctrica en el hospital, por lo que la asignación queda incompleta (el médico tiene asignado el quirófano pero el sistema no registró qué cirugía va a realizar). En este caso, el sistema debería tener una forma de volver al estado anterior a la asignación para no crear incongruencias; en otro caso, la cirugía tendría que ser aplazada o incluso se podría perder la información del turno, lo cual atenta contra la salud del paciente.

    \item \textbf{Anomalías de acceso concurrente (Concurrent-access anomalies):} Este punto estipula que en cualquier sistema la concurrencia de usuarios (el uso compartido del sistema por múltiples usuarios al mismo tiempo) es muy común y va cada vez más a la alza. Por ello, el sistema de un hospital debería tener la capacidad de que sus empleados puedan acceder a los datos de un paciente o trabajador al mismo tiempo sin que se generen incongruencias. Por ejemplo, si el departamento administrativo del hospital y una especialidad médica necesitan modificar la información de un paciente al mismo tiempo (digamos que la administración actualiza su número telefónico y el médico registra una nueva alergia detectada), sin un control de concurrencia podría pasar que ambos guarden su cambio sobre el mismo estado anterior, sobrescribiendo la información de la primera modificación guardada. En el peor de los casos, se perdería la alergia que el doctor acaba de diagnosticar sobre el paciente.

    \item \textbf{Problemas de seguridad (Security problems):} Este podría ser el punto más intuitivo u obvio, pero no por ello es menos importante. El sistema debe regular qué personal puede acceder a qué tipo de datos. Por ejemplo, el doctor puede acceder a los datos médicos generales del paciente, junto a sus tratamientos, citas y médicos asociados, sin embargo, no debería poder modificar sus datos personales administrativos como su número telefónico, correo o teléfonos de emergencia. Pero ya que la aplicación sigue un anticuado sistema de archivos muy específico para cada problema, cualquier restricción de seguridad que se intente aplicar a nivel hospitalario sería extremadamente complicada de implementar de forma homogénea.
\end{enumerate}
